
\section{Experimental Setup}
\label{sec:experimental-setup}

\subsection{Domains}
\label{sec:domains}

We evaluate on five benchmark domains spanning infrastructure operations and e-commerce, plus a production-scale domain for external validation (Table~\ref{tab:domain-stats}).

\begin{table}[t]
\centering
\small
\begin{tabular}{l rrr r}
\toprule
\textbf{Domain} & \textbf{Tools} & \textbf{Types} & \textbf{T/T\%} & \textbf{Pruning} \\
\midrule
Kubernetes & 135 & 41 & 30\% & 63\% \\
Ansible & 108 & 40 & 37\% & 67\% \\
GitHub & 133 & 44 & 33\% & 56\% \\
CI/CD & 54 & 51 & 94\% & 87\% \\
Shopify & 170 & 61 & 35\% & 78\% \\
\midrule
AAP MCP & 1,060 & 79 & 7\% & -- \\
\bottomrule
\end{tabular}
\caption{Domain statistics. T/T\% = types-to-tools ratio. Pruning = average entity pruning. CI/CD has near-parity types/tools (94\%) yet achieves the highest pruning (87\%), demonstrating that TCR derives advantage from graph topology rather than label-space reduction. AAP MCP is evaluated separately in Section~\ref{sec:production-eval}.}
\label{tab:domain-stats}
\end{table}

Benchmark registries are constructed from API documentation; Shopify (Admin REST API) serves as a non-infrastructure domain.
The \textbf{AAP MCP} domain is the official Ansible Automation Platform MCP Server\footnote{\url{https://github.com/ansible/aap-mcp-server}} (1,060 operations across four services), with its composition graph derived automatically from OpenAPI specifications.

\subsection{Models}
\label{sec:models}

We evaluate four language models spanning different architectures and sizes:
Granite 4.1 8B,
Qwen3-14B,
GPT-OSS-20B,
and Claude Haiku 4.5.
All models are evaluated off-the-shelf without fine-tuning.
Since TCR requires only entity type prediction, the approach is applicable to any model capable of natural language understanding, regardless of tool-use training.

\subsection{Benchmark Queries}
\label{sec:queries}

The benchmark contains 140 manually authored queries for the five benchmark domains (25--30 per domain) and 47 queries for the AAP domain.
Each query is annotated with ground-truth source and target entity types and the expected tool sequence.
Queries span six categories: \textbf{clean} (explicit requests), \textbf{multi-hop} (3+ tool compositions), \textbf{noisy} (distracting information), \textbf{synonym} (alternative terminology), \textbf{ambiguous} (multiple valid interpretations), and \textbf{multi-path} (multiple valid compositions).
All annotations were finalized before model evaluation began.

\subsection{Routing Strategies}
\label{sec:strategies}

We compare TCR against representative baselines:

\begin{itemize}
    \item \textbf{Baseline}: direct tool selection from the complete catalog via prompt-based selection.
    \item \textbf{Function Calling}: tools presented as OpenAI-compatible function schemas via the \texttt{tools} API parameter.
    \item \textbf{Retrieval}: top-10 tools by embedding similarity (nomic-embed-text-v1.5) presented for selection.
    \item \textbf{TCR} (ours): entity type prediction followed by graph search.
    \item \textbf{Oracle}: graph search with ground-truth entity types (upper bound).
\end{itemize}

\subsection{Metrics}
\label{sec:metrics}

For end-to-end routing: precision, recall, and F1 over predicted and expected tool sets.
For structural validity: hallucinated tool invocations (tools not in the registry).
For entity type prediction: source accuracy, target accuracy, and exact match (both correct).
For graph structure: entity pruning (Eq.~\ref{eq:entity-pruning}).

\subsection{Implementation}
\label{sec:implementation}

All LLM calls use greedy decoding (temperature 0) with identical prompt engineering. Results are deterministic across runs.
